%% kapa: exam=1 round=21 year=2010 period=2 form=A subject=appraisal_law subject_ko=부동산관계법규 seq=3 number=43
%% source: Q-Net 감정평가사 기출문제 (cst00309) artlSeq=5024181 file=extracted/1008/5024181/2010년 제21회 감정평가사 1차시험 기출문제/2교시.hwp
%% conversion: kordoc 4.24.51 -> tools/gen1.py (auto); review=proofread-gpt-sol; flags=none
\begin{question}{43}
국토의 계획 및 이용에 관한 법령상 도시계획에 관한 설명으로 옳은 것은?
\begin{choices}
  \item 도시계획은 광역도시계획, 도시기본계획, 도시관리계획으로 구분된다.
  \item 도시계획은 특별시․광역시․시 또는 군의 관할 구역에서 수립되는 다른 법률에 따른 토지의 이용․개발 및 보전에 관한 계획의 기본이 된다.
  \item 지방자치단체는 도시계획을 수립함에 있어 국토해양부장관이 평가한 도시의 지속가능성 평가결과를 반영하지 않아도 된다.
  \item 광역시장은 관할 구역에 대하여 다른 법률에 따른 환경에 관한 부문별 계획을 도시기본계획의 내용과 부합되지 않게 수립할 수 있다.
  \item 지식경제부장관이 20만 제곱미터의 용도지역의 변경에 관한 도시관리계획의 결정을 의제하는 경제자유구역개발계획을 승인하고자 하는 경우에는 중앙도시계획위원회의 심의를 받아야 한다.
\end{choices}
\end{question}
\answer{2}
